The upper inequality is \(\mathrm{thm:jackson-factorial-upper}\).  By \(\mathrm{def:minimax-risk}\), \(\mathfrak R_{n,d,epsilon}\) is the infimum of worst-case risks over all measurable estimators, and therefore
\[
 \mathfrak R_{n,d,epsilon}
 \leq\sup_{\mathbb P\in\mathcal D_{d,epsilon}^{\mathrm{obs}}}
 E_{\mathbb P}[(\widehat V_{n,d}^{\mathrm{JF}}-\Psi(\mathbb P))^2].
\]
The lower inequality is \(\mathrm{thm:all-estimator-lower}\).  Concatenating the three inequalities proves the displayed frontier for every \(n\geq1,d\geq2\).  Finally, \(\mathrm{prop:identification-and-extension}\) proves that the observed margins of \(\mathcal P_{d,epsilon}\) fill all of \(\mathcal D_{d,epsilon}^{\mathrm{obs}}\), and \(\mathrm{prop:causal-optimal-value-corollary}\) proves \(V^\star(P)=\Psi\{\operatorname{Obs}(P)\}\).  Thus the causal statement is a corollary of equality of the observation experiments and targets, not part of the statistical target's definition.